\section{Related work in full}
\label{app:related}

\paragraph{Offline metrics versus closed-loop competence.}
In neighboring embodied domains, it is established that offline imitation metrics can fail to
identify good closed-loop policies. In driving, offline prediction error is not necessarily
correlated with driving quality, and two models with identical prediction error can differ widely
when deployed \citep{codevilla2018offline}. Open-loop metrics on a modern planning benchmark favor
degenerate planners that nearly ignore their input \citep{dauner2023parting}. In manipulation, the
robomimic study reports that the best validation loss does not correspond to the best-performing
policy, with validation-selected checkpoints $10$--$100\%$ worse than rollout-selected ones
\citep{mandlekar2021robomimic}. \citet{hao2023masked} report that choosing a trained policy's
input-modality mask by lowest validation loss picked a zero-success mask in all three seeds on a
robomimic task, while choosing among separately trained policies by the same loss roughly doubled
another baseline's success. \citet{simchowitz2025pitfalls} construct a setting in which rollout cost
rises while validation loss falls throughout training. By 2026, a weak correlation between
action-space validation error and success is described as a practitioner consensus
\citep{huang2026critical}, although the same study finds raw validation error rank-consistent with
rollout success when only the number of training steps varies. Behavior-cloned policies can reach
near-zero held-out error while relying on spurious correlates that collapse under their own
deployment distribution \citep{dehaan2019causal}, including collapse toward a dominant action
regime overrepresented in demonstrations \citep{codevilla2019exploring}.
\citet{codevilla2018offline} also show that the offline--online correlation improves substantially
with better-chosen validation data and metrics. Our claim therefore concerns only the plain
cross-entropy over decisions in their natural action frequencies, which general-purpose fine-tuning
stacks compute when rollouts are unavailable. The divergence also has limits. Under log-loss with a
discrete action vocabulary, compounding error is at most polynomial in the horizon
\citep{foster2024behavior}, and the exponential-in-horizon bound of \citet{simchowitz2025pitfalls}
applies to continuous actions only. The polynomial bound does not rule out compounding error, or
the shift from expert to policy histories at deployment, and our experiments do not separate their
contribution from action imbalance. The mechanism we point to, that a frequency-dominated
cross-entropy is blind to a rare class, was stated for \textsc{stop} in the training loss by
\citet{xiang2020l2stop} and for rare action changes in validation accuracy by
\citet{wijmans2019eqa}, who also chose among five low-loss checkpoints by executing them and found
inflection-weighted cross-entropy a better predictor of end-to-end performance than plain
cross-entropy. None of this prior work on offline metrics combines instruction-following navigation,
a VLM policy and a success criterion that depends on a single rare action, which is the combination
studied here.

\paragraph{The stopping problem in VLN.}
Stopping failure is also documented within VLN, on the discrete navigation graph.
\citet{zhao2023mindthegap} quantify a persistent gap between Oracle Success and Success (agents
pass the goal and fail to stop) across four systems. Two further findings concern learned action
regularities. A case study of one agent finds its errors reflect a systematic bias toward forward
actions learned during training \citep{yang2023behavioral}, and \citet{thomason2019shifting} note
that peaked action distributions in R2R training let agents memorize simple rules, such as moving
forward an average number of steps. \citet{xiang2020l2stop} treat stopping as a distinct learning
problem and add a dedicated head. For a VLM fine-tuned for R2R-CE, \citet{dang2026rynnbrain} report
Oracle Success of $71.6\%$ against success of $58.6\%$ on \texttt{val\_unseen} and conclude that the
model ``lacks the precision for the terminal stopping maneuver''. They evaluate a trained model and
do not study how its checkpoint was chosen. Fine-tuning a small VLM for aerial navigation,
\citet{tamazyan2025aerial} find performance ``critically affected by a severe action imbalance in
the training data''. In VLM fine-tuning for R2R-CE, \citet{cheng2025navila} name the
underrepresentation of the stop action as a label-imbalance problem and report that removing their
data rebalancing lowers \texttt{val\_unseen} success from $49.7$ to $30.0$. Their remedy acts on the
data blend. The failure reported here, never stopping at the goal, survives inverse-frequency
\emph{loss} weighting (the weighted checkpoints press \textsc{stop} seven times, all outside the goal
radius), and the object of our study is the selection metric. Poor stopping is already documented, and we add a closed-loop
measurement, in continuous R2R-CE with a fine-tuned VLM and every saved
checkpoint executed, of what an averaged selection loss reports while the policies it ranks lose the
terminal action. The loss fell from $5.49$ nats to below $1$ in every run, and no checkpoint was
shown to navigate better than the one it chose. Every trained unweighted checkpoint predicts
\textsc{stop} on none of the $256$ selection decisions (update $0$ predicts it $28$ times), so in
those runs any rule choosing among trained checkpoints picks one that never predicts \textsc{stop}
there. The collapse is a training outcome, and the averaged loss did not flag it.

\paragraph{VLM policies for continuous navigation.}
Recent systems fine-tune large vision--language models to act directly in continuous environments,
building on Room-to-Room (R2R) and R2R-CE \citep{anderson2018r2r,krantz2020beyond,hong2022bridging}.
These systems differ widely in action interface, observation history and planning scaffolding, so we
make no comparison to published success rates and report only within-protocol contrasts against
references frozen in advance. R2R-CE papers that state a selection rule select on closed-loop
\texttt{val\_unseen} performance, and VLN-CE performs early stopping on it
\citep{krantz2020beyond}. NaVILA, a VLM policy, trains for one epoch and states no selection rule
\citep{cheng2025navila}. We found no R2R-CE paper that selects by teacher-forced loss. We study that
rule because it is what a general-purpose fine-tuning stack computes when rollouts are unavailable.
The Hugging Face Trainer, asked to keep its best checkpoint, keeps the one with the lowest
evaluation loss unless given another metric \citep{hf2026trainer}. Because \texttt{val\_unseen}
stays sealed, the closed-loop evaluation on reserved trajectories is a diagnostic of the rule, and
we make no leaderboard claim. Class-imbalance remedies such as re-weighting are standard
\citep{cui2019classbalanced}, and our second treatment arm applies one and measures what it changes.
The VLN-CE baseline itself trains with \emph{inflection weighting} \citep{wijmans2019eqa}, which
upweights timesteps where the action changes. The two remedies are not interchangeable. On our
$3{,}072$-decision training set every \textsc{stop} is an inflection, so inflection weighting would
give it $2.48\times$ the weight of a non-inflection step ($2.55\times$ over the released split, and
the VLN-CE baseline configures $3.2$ for R2R \citep{krantz2020beyond}), the same factor as any turn
onset. The inverse-frequency weighting we test gives \textsc{stop} $28.5\times$ the weight of
\textsc{move\_forward} on that same set. The two rules weight different groups of decisions, rare
classes against changes of action, so these factors do not rank the remedies' strength. We did not
test inflection weighting. The failure to stop at the goal survives the weighting we did test.
